\documentclass{article}
\usepackage[T1]{fontenc}
\usepackage{lmodern}
\usepackage{graphicx}
\usepackage{amsmath,amssymb}
\usepackage[a4paper,margin=2cm]{geometry}
\usepackage[absolute,overlay]{textpos}
\setlength{\TPHorizModule}{1mm}
\setlength{\TPVertModule}{1mm}
\usepackage[indonesian]{babel}
\usepackage{hyperref}
\setlength{\emergencystretch}{2em}
% unit: Halaman 92

\begin{document}

% === Halaman 92 (python/native) ===

• Jika pasangan plainteks dan cipherteks tidak tersedia, maka Affine 

Cipher masih bisa diserang dengan exhaustive key search. 

• sebab ada 25 pilihan untuk b dan 12 buah nilai m yang relatif prima 

dengan 26 (yaitu 1, 3, 5,  7, 9, 11, 15, 17, 19, 21, 23, dan 25). 

• Dengan mencoba sebanyak 25 x 12 = 300 susunan b dan m, maka 

nilai b dan m yang benar dapat ditemukan dengan mudah, karena 300 

susunan b dan m adalah jumlah yang sangat sedikit.

92

\end{document}
